\input{preamble.tex}

\usepackage{tikz}
\usetikzlibrary{positioning}
\definecolor {processblue}{cmyk}{0.96,0,0,0}

\begin{document}

% Indicate your name and your student number here (e.g., A01xxx)
\header{Assignment 1}{Wira Azmoon Ahmad}{A0149286R}

%======================================
\section{Shortest Path}
\subsection{Backward Dynamic Programming}
\begin{center}
\begin{tabular}{c c c c c c c c c} 
    \hline
      & A & B & C & D & E & F & G & H \\ [0.5ex] 
    \hline
    $V_0$ & 0 & $\infty$ & $\infty$ & $\infty$ & $\infty$ & $\infty$ 
    & $\infty$ & $\infty$ \\
    $V_1$ & 0 & $\infty$ & 7 & 7 & 5 & $\infty$ & $\infty$ & $\infty$ \\
    $V_2$ & 0 & 9 & 6 & 7 & 5 & 8 & 11 & 8 \\
    $V^*$ & 0 & 8 & 6 & 7 & 5 & 8 & 11 & 8 \\
    \hline
\end{tabular}
\end{center}

\subsection{Dijkstra's}

\begin {center}
\begin {tikzpicture}[auto ,node distance =4 cm and 5cm ,on grid ,
semithick ,
state/.style ={ circle , draw, minimum width =1 cm}]
\node[state] (A) {A};
\node[state] (E) [below right=of A] {E};
\node[state] (D) [below left=of A] {D};
\node[state] (B) [below left=3cm and 2cm of A] {B};
\node[state] (C) [below right=3cm and 2cm of A] {C};
\node[state] (F) [below=of B] {F};
\node[state] (G) [below=of C] {G};
\node[state] (H) [below=10cm of A] {H};
\path[]
    (A) edge node[left] {$7$} (D)
        edge node[right] {$5$} (E)
    (D) edge node[right] {$1$} (F)
        edge [bend right] node[right] {$1$} (H)
    (E) edge node[left] {$1$} (C)
        edge node[left] {$6$} (G)
    (C) edge node[left] {$2$} (B);
\end{tikzpicture}
\end{center}

%======================================
\newpage
\section{Quaternion}
\subsection{}
With $\theta=\frac{\pi}{2},(n_x,n_y,n_z)=\frac{1}{\sqrt{3}}(1,1,1)$,
\begin{equation}
\begin{aligned}
    h&=(\cos\frac{\theta}{2},n_x \sin\frac{\theta}{2},
    n_y \sin\frac{\theta}{2},n_z \sin\frac{\theta}{2})\\
    &=(\cos\frac{\pi}{4}, \frac{1}{\sqrt{3}}\sin\frac{\pi}{4},
    \frac{1}{\sqrt{3}}\sin\frac{\pi}{4},
    \frac{1}{\sqrt{3}}\sin\frac{\pi}{4})\\
    &=(\frac{\sqrt{2}}{2}, \frac{\sqrt{6}}{6}, 
    \frac{\sqrt{6}}{6}, \frac{\sqrt{6}}{6})
\end{aligned}
\end{equation}

\subsection{}
Yes.
\[h=(-\frac{\sqrt{2}}{2}, -\frac{\sqrt{6}}{6}, 
-\frac{\sqrt{6}}{6}, -\frac{\sqrt{6}}{6})\]

\subsection{}
\[h=(\frac{\sqrt{2}}{2}, -\frac{\sqrt{6}}{6}, 
-\frac{\sqrt{6}}{6}, -\frac{\sqrt{6}}{6})\]

%======================================
\newpage
\section{Configuration Space}
\subsection{}
\[\dim\mathcal{C}=2\]
The configuration can be specified by 
\[q=(\theta_1,\theta_2)\]
for the 2 revolute joints.
\subsection{}
\[\dim\mathcal{C}=12\]
Each robot has a configuration space of $\mathbb{R}^3\times SO(3)$,
which has dimension 6. Both together will have dimension 12.


\subsection{}
\[\dim\mathcal{C}=12\]
Each manipulator's configuration can be specified by
\[q=(\theta_1,\theta_2,\theta_3,\theta_4,\theta_5,\theta_6)\]
for 6 revolute joints, with dimension 6. 
Both together will have dimension 12.

%======================================
\newpage
\section{Goal Configuration}
The \textit{C-space goals} are denoted by the solid borders.
\subsection{}

\subsection{}

%======================================
\newpage
\section{Sample Points}
\subsection{Circle}

\subsection{Disc}

\subsection{N-Dimensional Disc}


%======================================
\newpage
\section{Rejection Sampling}


%======================================
\newpage
\section{PRM}


%======================================
\newpage
\section{Motion Planning Algorithms}
\subsection{}
Dynamic Programming can be used here.
The environment of the workshop should be relatively static,
and the C-space of the robot movement and handling should be small,
so it should not be computationally heavy.
The algorithm would be complete and simple to implement as well.

\subsection{}


%======================================
\newpage
\section{Hybrid A*}


%======================================
\newpage
\section{Heuristic}

\end{document}